\documentclass[12pt]{article}
\usepackage[T1]{fontenc}
\usepackage[utf8]{inputenc}
\usepackage{lmodern}
\usepackage{textcomp}
\usepackage{enumitem}
\usepackage{geometry}
\geometry{margin=1in}
\begin{document}
\begin{center}
Answer Key - Constitutional Law - Graduate Level Assessment 18 \
\small (Answers are ordered exactly as in the question paper.)
\end{center}

\section*{Multiple Choice}
\begin{enumerate}[label=\arabic*.]
\item \textbf{Answer:} B
\\ \textit{Options:} A) Be made by a merchant to a merchant.; B) Be contained in a signed writing which gives assurance that the offer will be held open.; C) State the period of time for which it is irrevocable.; D) Not be contained in a form supplied by the offeror.
\\ \textit{Note:} The correct answer is based on the principle under the Uniform Commercial Code that an irrevocable offer must be supported by a signed writing that explicitly assures the offeree that the offer will remain open for acceptance, thus providing legal certainty and enforceability.
\item \textbf{Answer:} C
\\ \textit{Options:} A) A law that is discriminatory on its face.; B) A discriminatory application of a facially neutral law.; C) A discriminatory effect.; D) A discriminatory motive.
\\ \textit{Note:} A discriminatory effect alone does not establish the requisite intent to discriminate under constitutional law, as intent requires showing that the government acted with a purpose to discriminate, rather than simply resulting in a discriminatory outcome.
\item \textbf{Answer:} B
\\ \textit{Options:} A) Death or insanity of a party.; B) Partial performance; C) Destruction of the subject matter of contract.; D) Supervening illegality
\\ \textit{Note:} Partial performance does not terminate a contract by operation of law because it indicates that at least some obligations under the contract have been fulfilled, thereby preserving the contract rather than nullifying it.
\item \textbf{Answer:} D
\\ \textit{Options:} A) Pay just compensation for the property; B) Terminate the regulation and pay damages that occurred while regulation was in effect; C) Pay double damages if not addressed in 30 days.; D) a or b
\\ \textit{Note:} The correct answer is "a or b" because, under the Fifth Amendment, if a regulation constitutes a taking of private property, the government is required to either provide just compensation to the property owner or demonstrate a legitimate public purpose for the taking.
\item \textbf{Answer:} B
\\ \textit{Options:} A) A common carrier delivering a new computer.; B) Crops and timber to be severed from the property next summer.; C) The purchase of a commercial property building.; D) The sale of an intangible asset.
\\ \textit{Note:} The correct answer is based on the UCC's application to contracts involving the sale of goods, and since crops and timber are considered goods that can be severed from the land, they fall under the UCC's provisions.
\end{enumerate}

\section*{True / False}
\begin{enumerate}[label=\arabic*.]
\setcounter{enumi}{5}
\item \textbf{Answer:} True
\\ \textit{Options:} A) True; B) False
\\ \textit{Note:} A bill of attainder is considered a legislative act that inflicts punishment without a trial because it specifically targets individuals or groups for punishment without the judicial process, violating the principle of due process enshrined in the Constitution.
\item \textbf{Answer:} True
\\ \textit{Options:} A) True; B) False
\\ \textit{Note:} The hierarchy of U.S. law is established by the Supremacy Clause of the Constitution, which dictates that the Constitution is the supreme law of the land, followed by federal laws and treaties, then executive agreements, and finally state laws, thus making the statement true.
\item \textbf{Answer:} False
\\ \textit{Options:} A) True; B) False
\\ \textit{Note:} The statement is false because for a dying declaration to be admissible, it must be made by a declarant who believes they are about to die, which is not established in the assertion that the testimony is merely hearsay.
\item \textbf{Answer:} False
\\ \textit{Options:} A) True; B) False
\\ \textit{Note:} The statement is false because the termination of a parent's custody rights and welfare benefits typically requires a showing of clear and convincing evidence of parental unfitness, as well as compliance with due process requirements beyond mere notice and opportunity to respond.
\item \textbf{Answer:} True
\\ \textit{Options:} A) True; B) False
\\ \textit{Note:} Solicitation is considered a crime because it involves the act of enticing or encouraging another person to commit a crime, regardless of whether an actual agreement or conspiracy exists between multiple parties.
\end{enumerate}

\section*{Long Form Response}
\begin{enumerate}[label=\arabic*.]
\setcounter{enumi}{10}
\item \textbf{Answer:} In negligence cases, a defendant's unique circumstances, such as being deaf, significantly affect the standard of care expected of them. The law recognizes that individuals with disabilities may face different challenges and thus should be assessed based on what a reasonable person with similar conditions would do. For a deaf defendant accused of failing to heed a warning bell, the court would consider whether there were alternative means of awareness available and whether the defendant took reasonable steps to mitigate the risk of harm. This approach aligns with constitutional principles of fairness and equal protection, ensuring that individuals are not held to a standard that fails to account for their disabilities. Ultimately, the standard of care must be contextualized within the defendant's capabilities, reinforcing the need for a nuanced application of negligence law that respects individual differences.
\\ \textit{Note:} correct because it emphasizes that the standard of care in negligence cases is modified to account for a defendant's unique circumstances, such as being deaf, by evaluating their actions against the reasonable conduct expected from someone with similar disabilities, thereby aligning with principles of fairness and justice in constitutional law.
\item \textbf{Answer:} Judicial notice is a legal doctrine that allows a court to recognize certain facts as true without requiring formal proof, typically encompassing facts that are universally known, readily verifiable, or capable of accurate and ready determination. Examples of appropriate facts for judicial notice include historical events, geographic facts, and scientific principles. In contrast, facts asserted by individual political organizations do not qualify for judicial notice due to their inherent subjectivity and potential bias; such assertions may lack the objectivity and verifiability required for judicial notice, which aims to maintain impartiality in legal proceedings. Therefore, while facts recognized by political entities may carry significance in a political context, they do not meet the stringent criteria necessary for judicial notice in a trial.
\\ \textit{Note:} correct because it accurately defines judicial notice as a mechanism for recognizing universally accepted or easily verifiable facts while explaining that assertions from individual political organizations lack the objective certainty and broad consensus required for judicial notice.
\end{enumerate}

\vfill
\noindent\textit{End of Answer Key}
\end{document}